\chapter{Pipeline dan Pipe-and-Filter}
\label{bab:pipa}

\section*{Tujuan Pembelajaran}

\begin{enumerate}
  \item Menjelaskan pembagian antara \textit{filter} dan \textit{pipe}
    beserta \textit{contract} yang menyambung keduanya.
  \item Membedakan pemrosesan \textit{streaming} dari pemrosesan \textit{batch},
    lalu menentukan kapan masing-masing layak dipilih.
  \item Mengukur \textit{filter independence}, akibat \textit{stage reordering},
    dan \textit{first output latency} pada kedua cara pemrosesan.
\end{enumerate}

\section{Pendahuluan}

\textit{Pipe-and-filter} menata sebuah pekerjaan menjadi deretan \textit{stage}
yang masing-masing mengerjakan satu perubahan kecil. Setiap \textit{stage} disebut
\textit{filter}, yaitu penyaring satu arah, sedangkan penyambung antar
\textit{stage} disebut \textit{pipe}. Rangkaian keduanya disebut
\textit{pipeline}. Data mengalir dari \textit{stage} pertama sampai
\textit{stage} terakhir, dan tidak ada \textit{stage} yang memanggil \textit{stage}
lain secara langsung.

Bentuk ini sudah berumur lebih dari lima puluh tahun. Ritchie mencatat bahwa
\textit{pipeline} muncul di Unix pada tahun 1972 atas usulan McIlroy, yang lama menganjurkan
alur kendali tak berjenjang seperti pada
korutin~\cite{ritchie1984evolution}. Gagasan yang sama tetap dipakai sampai
sekarang. Pustaka pembelajaran mesin scikit-learn, misalnya, menyediakan kelas
\texttt{Pipeline} tersendiri, dan menuntut setiap \textit{stage} selain
\textit{stage} terakhir memiliki metode
\texttt{transform}~\cite{sklearn2026pipeline}.

Bab ini memakai contoh yang dekat dengan pekerjaan sehari-hari, yaitu
penyeragaman \textit{file} gambar. Foto datang dari sumber yang beragam, lalu
diubah menjadi beberapa resolusi baku dengan tanda air. Bahannya enam puluh foto
yang diunduh \texttt{unduh\_gambar.py} dari layanan Lorem Picsum memakai
\textit{seed} tetap, yaitu \texttt{apl1} sampai \texttt{apl60} pada ukuran
1600 kali 1200 piksel, sehingga kumpulan gambarnya sama pada setiap mesin dan
angka latihannya dapat dibandingkan antar mahasiswa. Pembahasan mengikuti
kerangka yang disusun oleh Richards dan Ford~\cite{richards2020fundamentals}.

\section{\textit{Filter} dan \textit{Pipe}}

Pembagiannya hanya dua. \textbf{\textit{Filter}} menerima satu
\textit{stream} lalu menghasilkan \textit{stream} baru, dan tidak mengetahui siapa
yang mengirim maupun siapa yang menerima. \textbf{\textit{Pipe}} menyambungkan
keluaran satu \textit{filter} menjadi masukan \textit{filter} berikutnya, dan
tidak mengubah isinya sama sekali.

Gambar~\ref{fig:pipa-lurus} menunjukkan susunan \textit{pipeline} baku pada contoh bab ini.
Panahnya searah, dan tidak ada panah yang kembali ke belakang.

\begin{figure}[htbp]
  \centering
  \begin{tikzpicture}[
    tahap/.style={draw, rounded corners=2pt, minimum width=2.2cm,
      minimum height=1.1cm, align=center, font=\scriptsize\bfseries},
    panah/.style={-{Stealth[length=2mm]}, thick}
  ]
    \node[tahap, fill=gray!22] (sumber) {masukan/\\
      {\scriptsize\mdseries 60 file}};
    \node[tahap, fill=blue!8, right=0.3cm of sumber] (mode) {ubah\_mode\\
      {\scriptsize\mdseries ke RGB}};
    \node[tahap, fill=blue!8, right=0.3cm of mode] (kecil) {Resize\\
      {\scriptsize\mdseries lebar 800}};
    \node[tahap, fill=blue!8, right=0.3cm of kecil] (tanda) {Watermark\\
      {\scriptsize\mdseries menyatukan}};
    \node[tahap, fill=green!10, right=0.3cm of tanda] (simpan) {SaveJpeg\\
      {\scriptsize\mdseries kualitas 85}};

    \draw[panah] (sumber) -- (mode);
    \draw[panah] (mode) -- (kecil);
    \draw[panah] (kecil) -- (tanda);
    \draw[panah] (tanda) -- (simpan);
  \end{tikzpicture}
  \caption{\textit{Pipeline} baku pada contoh bab ini. Setiap \textit{stage} hanya mengenal bentuk
    data yang lewat, bukan \textit{stage} sebelum maupun sesudahnya}
  \label{fig:pipa-lurus}
\end{figure}

Analogi yang mendekati adalah lini pengolahan buah. Buah datang di atas
\textit{conveyor belt}, lalu melewati meja pemilah yang memisahkan jeruk dari
tomat. Jeruk diarahkan ke mesin pemeras menjadi jus, sedangkan tomat diarahkan
ke mesin penghalus menjadi saus. Setiap mesin hanya menangani buah yang datang
kepadanya, dan tidak mengetahui proses di mesin yang lain.
\textit{Conveyor belt}-nya sendiri tidak memproses buah apa pun.

Pembagian tersebut menentukan letak setiap keputusan.
Tabel~\ref{tab:penyaring-pipa} merangkum tanggung jawab dan larangan kedua
bagian.

\begin{table}[htbp]
  \centering
  \caption{Tanggung jawab dan larangan \textit{filter} serta \textit{pipe}}
  \label{tab:penyaring-pipa}
  \begin{tabularx}{\textwidth}{|l|X|X|}
    \hline
    \textbf{Bagian} & \textbf{Tanggung jawab} & \textbf{Tidak boleh mengetahui} \\
    \hline
    \textit{Filter} & Satu perubahan atas data yang lewat & Nama \textit{stage} lain, dan letaknya di dalam susunan \\
    \hline
    \textit{Pipe} & Menyambung keluaran menjadi masukan berikutnya & Isi perubahan yang dikerjakan setiap \textit{stage} \\
    \hline
    \textit{Contract} & Bentuk data yang mengalir di antara \textit{stage} & Teknologi penyimpanan maupun sumber datanya \\
    \hline
  \end{tabularx}
\end{table}

\section{\textit{Contract} yang Menyambung}

\textit{Contract}-nya hanya satu kalimat, yaitu setiap \textit{filter}
menerima satu \textit{iterable} lalu menghasilkan satu \textit{iterable}.
\textit{Filter} tanpa pengaturan ditulis sebagai fungsi generator, sedangkan
\textit{filter} yang membutuhkan pengaturan ditulis sebagai kelas yang dapat
dipanggil. \textit{Pipe} memperlakukan keduanya dengan cara yang sama persis.

Listing~\ref{lst:penyaring} memperlihatkan dua bentuk tersebut berdampingan.

\begin{lstlisting}[style=PythonStyle, caption={Dua bentuk \textit{filter},
  tersedia di \berkas{sources/chapter06/penyaring.py}}, label=lst:penyaring]
def ubah_mode(masukan):
  """Menyeragamkan mode warna setiap gambar menjadi RGB."""
  for nama, gambar in masukan:
    yield nama, gambar.convert(MODE_BAKU) if gambar.mode != MODE_BAKU else gambar


class Resize:
  """Filter yang menyeragamkan lebar gambar sambil menjaga rasionya."""

  def __init__(self, lebar: int) -> None:
    self.lebar = lebar

  def __call__(self, masukan):
    for nama, gambar in masukan:
      tinggi = round(gambar.height * self.lebar / gambar.width)
      yield nama, gambar.resize((self.lebar, tinggi), Image.LANCZOS)
\end{lstlisting}

Penyambungnya sependek itu pula. \textit{Pipe} hanya menumpuk pemanggilan, dan
penumpukan itu dikerjakan malas, sehingga belum satu \textit{file} pun dibuka ketika
penyusunan selesai.

\begin{lstlisting}[style=PythonStyle, caption={Penyambung seluruh \textit{stage},
  tersedia di \berkas{sources/chapter06/pipa.py}}, label=lst:sambung]
def jalankan(sumber: str, tahap: list, batas: int | None = None):
  """Menyambung seluruh stage menjadi satu pipeline, lalu mengembalikannya."""
  aliran = penyaring.baca_gambar(sumber, batas)
  for fungsi in tahap:
    aliran = fungsi(aliran)
  return aliran
\end{lstlisting}

Gambar~\ref{fig:kelas-pipa} menunjukkan hubungan keduanya sebagai
\textit{class diagram}. Perhatikan bahwa tidak ada satu pun panah antar
\textit{filter}.

\begin{figure}[htbp]
  \centering
  \includegraphics[width=0.92\textwidth, height=0.42\textheight,
    keepaspectratio]{kelas-pipa.pdf}
  \caption{\textit{Pipe} mengenal \textit{contract} pemanggilan saja. Kelima
    \textit{filter} tidak saling menunjuk, dan justru itulah yang membuat
    urutannya dapat diubah}
  \label{fig:kelas-pipa}
\end{figure}

Kemandirian tersebut dapat dihitung, bukan sekadar diyakini. Skrip
\texttt{periksa\_pipa.py} membaca kode tanpa menjalankannya, lalu melaporkan
\textit{filter} mana yang menyebut nama \textit{filter} lain. Pada contoh bab ini angkanya
nol dari lima, dan Latihan~1 mengulang pemeriksaan itu.

\section{Urutan \textit{Stage} adalah Keputusan}

\textit{Contract} yang seragam membuka satu kemungkinan yang menggoda, yaitu menukar
urutan \textit{stage}. Penukaran itu sah selama kedua \textit{stage} menerima dan menghasilkan
bentuk data yang sama. Pertanyaannya, apakah penukaran yang sah selalu aman.

Contoh bab ini memakai dua \textit{stage} yang sama-sama menerima dan menghasilkan
gambar, yaitu pengecilan dan penempelan tanda air. Keduanya boleh ditukar, dan
hasil penukarannya diukur pada Latihan~2. Hasil pengukuran nyatanya memuat dua
kejanggalan sekaligus, dan keduanya layak dibaca bersama-sama.

Kejanggalan pertama, selisih waktunya kecil sekali. Median urutan
\texttt{Resize} lebih dahulu tercatat 0,5835 detik, sedangkan urutan
\texttt{Watermark} lebih dahulu 0,6054 detik. Selisihnya 0,0219 detik, padahal
dua putaran pada urutan yang sama saja dapat berselisih 0,3660 detik. Selisih
antara putaran tercepat dan terlambat itulah yang disebut \textit{range}, dan
\textit{range} di sini 17 kali selisih mediannya.

Kejanggalan kedua justru menunjuk arah yang jelas. Pencacahan per putaran
mencatat urutan \texttt{Resize} lebih cepat pada 88 dari 100 putaran. Bila
kedua urutan sama cepatnya, peluang munculnya angka sejauh itu dari 50 kurang
dari dua per sepuluh pangkat lima belas, menurut uji tanda atas seratus putaran
berpasangan. Arah selisihnya karena itu bukan kebetulan.

Kedua ukuran tersebut dibaca bersama-sama. Arahnya tetap, yaitu urutan
\texttt{Resize} lebih dahulu memang lebih cepat, sedangkan besarnya kecil, yaitu
3,8 persen dari mediannya dan 6 persen dari \textit{range}-nya. Selisih sekecil
itu tidak akan dirasakan pengguna, sehingga urutan \textit{stage} di sini tidak
layak dipilih berdasarkan waktu.

Yang tidak meragukan justru hasilnya. \textit{File} keluaran kedua urutan tidak
sama bait demi bait, dan sidik keduanya terulang sama persis pada setiap
penjalanan. Sebabnya, menempelkan tanda air pada kanvas besar lalu mengecilkan
seluruhnya membuat tanda air itu ikut diperhalus. Gambar~\ref{fig:kotak-urutan}
memperlihatkan sebaran keduanya berdampingan.

\begin{figure}[htbp]
  \centering
  \resizebox{\textwidth}{!}{%
  \begin{tikzpicture}[
      kotak/.style={draw=praditagreen, thick},
      sumbu/.style={draw=gray!70},
      label/.style={font=\scriptsize},
      kecil/.style={font=\tiny\mdseries}
    ]
      % Sumbu waktu, 1cm mewakili 0,05 detik mulai dari 0,35 detik.
      \draw[sumbu] (0,0) -- (10,0);
      \foreach \x/\teks in {0/0{,}35, 2/0{,}45, 4/0{,}55, 6/0{,}65, 8/0{,}75,
        10/0{,}85} {
        \draw[sumbu] (\x,0) -- (\x,-0.12);
        \node[kecil, below] at (\x,-0.12) {\teks};
      }
      \node[kecil, below] at (5,-0.6) {waktu satu putaran (detik)};
  
      % Urutan Resize dahulu: min 0,3893  Q1 0,5625  median 0,5835
      % Q3 0,6104  maks 0,7553
      \draw[kotak] (0.786,1.75) -- (4.25,1.75);
      \draw[kotak] (5.208,1.75) -- (8.106,1.75);
      \draw[kotak] (0.786,1.6) -- (0.786,1.9);
      \draw[kotak] (8.106,1.6) -- (8.106,1.9);
      \draw[kotak, fill=blue!8] (4.25,1.5) rectangle (5.208,2.0);
      \draw[kotak, very thick] (4.67,1.5) -- (4.67,2.0);
      \node[label, left] at (-0.15,1.75) {\texttt{Resize} dahulu};
  
      % Urutan Watermark dahulu: min 0,4065  Q1 0,5827  median 0,6054
      % Q3 0,6480  maks 0,8049
      \draw[kotak] (1.13,0.75) -- (4.654,0.75);
      \draw[kotak] (5.96,0.75) -- (9.098,0.75);
      \draw[kotak] (1.13,0.6) -- (1.13,0.9);
      \draw[kotak] (9.098,0.6) -- (9.098,0.9);
      \draw[kotak, fill=green!10] (4.654,0.5) rectangle (5.96,1.0);
      \draw[kotak, very thick] (5.108,0.5) -- (5.108,1.0);
      \node[label, left] at (-0.15,0.75) {\texttt{Watermark} dahulu};
    \end{tikzpicture}}
  \caption{Sebaran waktu seratus putaran kedua urutan. Kotaknya adalah kuartil pertama sampai ketiga, garis tebal di dalamnya median, dan kedua ujung kumis adalah nilai terkecil serta terbesar}
  \label{fig:kotak-urutan}
\end{figure}

\section{\textit{Fan-out} dan \textit{Merge}}

\textit{Pipeline} yang benar-benar lurus jarang bertahan lama. Satu \textit{file} sumber biasanya
dibutuhkan dalam beberapa ukuran sekaligus, misalnya kecil untuk daftar, sedang
untuk halaman, dan besar untuk tampilan penuh. Kebutuhan itu dijawab dengan
\textbf{\textit{fan-out}}, yaitu satu \textit{stream} digandakan menjadi
beberapa \textit{stream} yang berjalan sendiri-sendiri.

Lawannya adalah \textbf{\textit{merge}}, yaitu beberapa \textit{stream}
disatukan menjadi satu. Pada contoh bab ini \textit{merge} terjadi dua kali. Yang
pertama di dalam setiap cabang, yaitu ketika foto disatukan dengan \textit{file}
tanda air lewat \textit{alpha channel}. Yang kedua di ujung, yaitu ketika
laporan ketiga
cabang disambung menjadi satu \textit{stream}.

Gambar~\ref{fig:cabang-tiga} memperlihatkan susunannya sebagai blok.

\begin{figure}[htbp]
  \centering
  \resizebox{\textwidth}{!}{%
  \begin{tikzpicture}[
      kotak/.style={draw, rounded corners=2pt, minimum width=2.1cm,
        minimum height=0.85cm, align=center, font=\scriptsize},
      hulu/.style={kotak, fill=blue!8},
      cabang/.style={kotak, fill=blue!8, minimum width=3.1cm},
      simpul/.style={draw, circle, fill=praditagreen!25, minimum size=0.85cm,
        inner sep=0pt, font=\scriptsize\bfseries},
      hasil/.style={kotak, fill=green!10, minimum width=2.4cm},
      panah/.style={-{Stealth[length=2mm]}, thick}
    ]
      \node[hulu] (sumber) {masukan/\\{\tiny\mdseries 60 file}};
      \node[hulu, right=0.5cm of sumber] (mode) {ubah\_mode\\{\tiny\mdseries ke RGB}};
      \node[simpul, right=0.5cm of mode] (tee) {tee};
  
      \node[cabang, right=0.7cm of tee] (sedang)
        {Resize(800), Watermark,\\SaveJpeg(85)};
      \node[cabang, above=0.35cm of sedang] (kecil)
        {Resize(320), Watermark,\\SaveJpeg(75)};
      \node[cabang, below=0.35cm of sedang] (besar)
        {Resize(1600), Watermark,\\SaveJpeg(90)};
  
      \node[simpul, right=0.7cm of sedang] (chain) {chain};
      \node[hasil, right=0.5cm of chain] (laporan)
        {laporan gabungan\\{\tiny\mdseries 180 file}};
  
      \draw[panah] (sumber) -- (mode);
      \draw[panah] (mode) -- (tee);
      \draw[panah] (tee) -- (kecil);
      \draw[panah] (tee) -- (sedang);
      \draw[panah] (tee) -- (besar);
      \draw[panah] (kecil) -- (chain);
      \draw[panah] (sedang) -- (chain);
      \draw[panah] (besar) -- (chain);
      \draw[panah] (chain) -- (laporan);
  
      \node[font=\tiny\mdseries, below=0.1cm of tee] {\textit{fan-out}};
      \node[font=\tiny\mdseries, below=0.1cm of chain] {\textit{merge}};
    \end{tikzpicture}}
  \caption{Satu sumber dibaca sekali, digandakan menjadi tiga cabang resolusi, lalu laporan ketiganya disatukan kembali}
  \label{fig:cabang-tiga}
\end{figure}

Gambar~\ref{fig:sekuens-percabangan} menunjukkan urutannya. \textit{Stage} penyeragaman
mode dikerjakan sekali di hulu, sebab hasilnya sama bagi ketiga cabang.

\begin{figure}[htbp]
  \centering
  \includegraphics[width=\textwidth, height=0.46\textheight,
    keepaspectratio]{sekuens-percabangan.pdf}
  \caption{Satu kali pembacaan sumber melayani tiga cabang resolusi, lalu
    laporan ketiganya disambung kembali menjadi satu \textit{stream}}
  \label{fig:sekuens-percabangan}
\end{figure}

Hasil nyatanya tercantum pada Tabel~\ref{tab:cabang}. Enam puluh \textit{file} sumber
menghasilkan seratus delapan puluh \textit{file} keluaran, dan sumbernya cukup dibaca
sekali.

\begin{table}[htbp]
  \centering
  \caption{Tiga cabang resolusi dari satu kali pembacaan sumber}
  \label{tab:cabang}
  \begin{tabularx}{\textwidth}{|l|l|l|X|}
    \hline
    \textbf{Cabang} & \textbf{Lebar} & \textbf{Kualitas} & \textbf{Ukuran seluruhnya} \\
    \hline
    kecil & 320 & 75 & 723 KiB \\
    \hline
    sedang & 800 & 85 & 5.118 KiB \\
    \hline
    besar & 1600 & 90 & 15.697 KiB \\
    \hline
  \end{tabularx}
\end{table}

\section{\textit{Streaming} atau \textit{Batch}}

Satu \textit{pipeline} yang sama dapat dijalankan dengan dua cara.
\textbf{\textit{Streaming}} berarti satu \textit{file} menempuh seluruh
\textit{stage} sebelum \textit{file} berikutnya dibuka. \textbf{\textit{Batch}}
berarti satu \textit{stage} dituntaskan untuk seluruh \textit{file}, baru
\textit{stage} berikutnya mulai. Keduanya memakai \textit{stage} yang sama
persis, dan hanya cara penyambungannya yang berbeda.

Perbedaannya tidak terlihat pada waktu total, sebab pekerjaannya memang sama.
Perbedaannya terlihat pada dua hal lain, yaitu kapan \textit{file} pertama selesai dan
berapa banyak yang harus ditampung di dalam memori sekaligus.

Kasus berikut memperlihatkan akibatnya. Sebuah tim menulis pengolah foto untuk
sistem penerimaan mahasiswa baru. Pengolahnya membaca seluruh \textit{file} unggahan
menjadi satu daftar, menyeragamkan ukurannya, menempelkan tanda air, lalu
menyimpannya. Pada pengujian dengan tiga puluh \textit{file} semuanya berjalan mulus.
Pada hari pendaftaran, \textit{file} yang masuk mencapai belasan ribu. Proses
berhenti karena kehabisan memori, dan tidak satu pun \textit{file} keluaran
tertulis, sebab \textit{stage} penyimpanan memang belum sempat dimulai. Andai
\textit{pipeline}-nya \textit{streaming}, \textit{file} pertama sudah tersimpan
sejak detik pertama.

Ukuran nyatanya tercantum pada Tabel~\ref{tab:aliran}. \textit{File} pertama
pada jalur \textit{streaming} selesai 9,5 kali lebih cepat, dan puncak
\textit{Resident Set Size} (RSS) prosesnya berbanding 118,4 melawan 576,8 MiB.

\begin{table}[htbp]
  \centering
  \caption{Pemrosesan \textit{streaming} dibanding \textit{batch}, seratus putaran}
  \label{tab:aliran}
  \begin{tabularx}{\textwidth}{|X|l|l|}
    \hline
    \textbf{Yang diukur} & \textbf{\textit{Streaming}} & \textbf{\textit{Batch}} \\
    \hline
    Median \textit{file} pertama & 83,1 ms & 792,8 ms \\
    \hline
    Puncak RSS proses & 118,4 MiB & 576,8 MiB \\
    \hline
    Putaran lebih cepat & 100 dari 100 & 0 dari 100 \\
    \hline
  \end{tabularx}
\end{table}

Angka memori tersebut diambil lewat \texttt{resource.getrusage}, yang
melaporkan puncak RSS proses, yaitu memori yang benar-benar dipegang proses
menurut sistem operasi. Setiap jalur diukur di dalam proses anak tersendiri,
sebab puncak RSS adalah tanda air tertinggi yang tidak pernah turun, sehingga
mengukur kedua jalur di dalam satu proses akan melaporkan angka jalur termahal
dua kali. \texttt{tracemalloc} sengaja tidak dipakai, sebab penyangga piksel
Pillow dialokasikan di luar jangkauannya lalu menyembunyikan selisih yang justru
menjadi pokok bahasan.

\section{Ringkasan}

\textit{Pipe-and-filter} menata pekerjaan menjadi deretan \textit{stage} yang
masing-masing mengerjakan satu perubahan kecil. \textit{Filter} tidak mengenal
tetangganya, dan \textit{pipe}-lah yang menyambung mereka. Gagasan itu berumur
lebih dari lima puluh tahun, lahir sebagai \textit{pipeline} Unix pada 1972, dan
masih dipakai pada perkakas masa kini.

Yang membuat susunan tersebut bekerja adalah \textit{contract} yang seragam,
yaitu satu \textit{stream} masuk dan satu \textit{stream} keluar. Karena
\textit{contract}-nya sama, \textit{stage} dapat ditambah, dibuang, atau ditukar
tanpa menyunting isi \textit{stage} lain. Latihan~1 menghitung
\textit{filter independence} itu, dan angkanya nol penyebutan dari lima
\textit{filter}.

Kebebasan itu punya batas yang sering terlewat. Dua \textit{stage} hanya boleh
ditukar bila bentuk data yang dilewatkannya sama, dan penukaran yang sah pun
belum tentu menghasilkan \textit{file} yang sama. Latihan~2 memperlihatkan
keduanya sekaligus, termasuk selisih waktu yang tenggelam di dalam
\textit{range} pengukurannya sendiri.

Cara penyambungan menentukan hal yang tidak tampak pada waktu total.
\textit{Streaming} menyerahkan \textit{file} pertama 9,5 kali lebih cepat, dan
puncak RSS prosesnya 118,4 MiB melawan 576,8 MiB pada jalur \textit{batch}.
\textit{Batch} baru masuk akal ketika satu \textit{stage} memang membutuhkan
seluruh data, misalnya pengurutan.

Pesan utama bab ini karena itu bukan bahwa \textit{pipeline} membuat program lebih cepat,
melainkan bahwa \textit{pipeline} memindahkan keputusan ke satu tempat. Urutan \textit{stage}, titik
percabangan, dan cara penyambungan seluruhnya ditentukan di luar \textit{filter}, dan
di situlah letak kelenturannya.

\section{Latihan}

Ketiga latihan berikut memakai satu basis kode yang sama di
\berkas{sources/chapter06/}, tetapi menyasar tiga masalah yang berbeda.
Latihan~1 memeriksa \textit{filter independence}. Latihan~2 memeriksa akibat
\textit{stage reordering}. Latihan~3 mengukur \textit{first output latency} pada
kedua cara pemrosesan. Hasil setiap mahasiswa dapat berbeda, dan yang dilaporkan adalah pola
perbandingannya, bukan angka mutlaknya.

Penyiapannya dijalankan sekali seperti pada Listing~\ref{lst:venv-ch06}. Bab ini
memakai Pillow, yang sudah tercantum pada \textit{file} kebutuhan.

\begin{lstlisting}[language=bash, caption={Menyiapkan bahan sebelum ketiga
  latihan, seluruh skrip tersedia di \berkas{sources/chapter06/}},
  label=lst:venv-ch06]
bash sources/siapkan_venv.sh
source .venv/bin/activate
cd sources/chapter06
python unduh_gambar.py
# Keluarannya: Gambar baru diunduh : 60
python buat_tanda_air.py
python pipa.py hitung
# Keluarannya baris terakhir: Keluaran akhir: 60 file
\end{lstlisting}

Bila jaringan tidak tersedia, jalankan \texttt{python buat\_gambar.py} sebagai
gantinya. Skrip itu menghasilkan gambar sintetis berukuran serupa, sehingga
ketiga latihan tetap dapat dikerjakan.

\subsection*{Latihan 1: \textit{Filter Independence}}

Latihan ini menguji klaim utama \textit{pipe-and-filter}, yaitu \textit{stage} dapat
ditambah dan ditukar tanpa menyunting \textit{stage} lain. Klaim itu gugur bila satu
\textit{filter} menyebut nama \textit{filter} lain, sebab penyebutan seperti itu mengunci
urutannya. Langkah kerjanya sebagai berikut.

\begin{enumerate}
  \item Jalankan kedua perintah pada Listing~\ref{lst:jalankan-periksa-pipa},
    lalu catat keempat angka yang dilaporkan.
  \item Buka \berkas{sources/chapter06/penyaring.py}, lalu periksa sendiri
    apakah ada \textit{filter} yang menyebut nama \textit{filter} lain.
  \item Tambahkan satu \textit{filter} baru yang memutar gambar, misalnya kelas
    \texttt{Rotate}, lalu sisipkan ke dalam susunan pada
    \berkas{sources/chapter06/pipa.py}.
  \item Jalankan ulang kedua perintah tersebut, lalu catat berapa \textit{file}
    \textit{filter} lain yang perlu disunting selama langkah ketiga.
\end{enumerate}

\begin{lstlisting}[language=bash, caption={Memeriksa \textit{filter
  independence}, skrip tersedia di \berkas{sources/chapter06/periksa_pipa.py}},
  label=lst:jalankan-periksa-pipa]
python periksa_pipa.py
# Keluarannya: Filter ditemukan      : 5
#              Filter menyebut filter: 0 {}
#              Baris filter          : 96
#              Baris pipe            : 109
python pipa.py hitung
# Keluarannya baris terakhir: Keluaran akhir: 60 file
\end{lstlisting}

Tabel~\ref{tab:lat1pipa-contoh} berisi hasil satu kali pemeriksaan nyata.

\begin{table}[htbp]
  \centering
  \caption{Contoh hasil Latihan 1 yang sudah terisi}
  \label{tab:lat1pipa-contoh}
  \begin{tabularx}{\textwidth}{|X|l|}
    \hline
    \textbf{Yang dilaporkan} & \textbf{Nilai} \\
    \hline
    \textit{Filter} ditemukan & 5 \\
    \hline
    \textit{Filter} menyebut \textit{filter} & 0 \\
    \hline
    Baris \textit{filter} & 96 \\
    \hline
    Baris \textit{pipe} & 109 \\
    \hline
  \end{tabularx}
\end{table}

Isikan hasil pemeriksaan pada Tabel~\ref{tab:lat1pipa-hasil}.

\begin{table}[htbp]
  \centering
  \caption{Lembar isian Latihan 1}
  \label{tab:lat1pipa-hasil}
  \begin{tabularx}{\textwidth}{|X|l|l|}
    \hline
    \textbf{Yang dilaporkan} & \textbf{Sebelum} & \textbf{Sesudah} \\
    \hline
    \textit{Filter} ditemukan & \dots & \dots \\
    \hline
    \textit{Filter} menyebut \textit{filter} & \dots & \dots \\
    \hline
    \textit{File} \textit{filter} disunting & \dots & \dots \\
    \hline
    \textit{File} \textit{pipe} disunting & \dots & \dots \\
    \hline
  \end{tabularx}
\end{table}

Jawab tiga pertanyaan berikut dengan menunjuk angka pada tabel. Pertama,
sebutkan bagian mana yang membuat penambahan \textit{stage} tidak menyentuh \textit{filter}
lain. Kedua, jelaskan apa yang terjadi pada angka penyebutan bila satu \textit{filter}
memanggil \textit{filter} lain di dalam tubuhnya demi kemudahan. Ketiga, bandingkan
jumlah baris \textit{filter} dengan jumlah baris \textit{pipeline}, lalu jelaskan mengapa
penyambungnya justru tidak lebih pendek.

\subsection*{Latihan 2: \textit{Stage Reordering}}

Latihan ini membuktikan bahwa penukaran yang sah menurut \textit{contract} belum tentu
menghasilkan \textit{file} yang sama. Kedua \textit{stage} yang ditukar menerima dan
menghasilkan bentuk data yang sama, sehingga \textit{pipeline} menerimanya tanpa keberatan.
Langkah kerjanya sebagai berikut.

\begin{enumerate}
  \item Jalankan perintah pada Listing~\ref{lst:jalankan-susun-ulang}. Skrip
    tersebut menjalankan seratus putaran untuk kedua urutan, dan mencetak
    kemajuan setiap dua puluh putaran.
  \item Catat median, \textit{range}, pencacahan putaran, dan kedua sidik
    \textit{file}-nya.
  \item Buka salah satu \textit{file} keluaran dari kedua direktori, lalu
    bandingkan ketajaman tanda airnya dengan mata.
  \item Ubah susunan pada \berkas{sources/chapter06/pipa.py} sehingga
    \texttt{SaveJpeg} berada sebelum \texttt{Watermark}, jalankan, lalu catat
    \textit{error} yang muncul. Satu kali percobaan nyata berhenti pada
    \texttt{ValueError: too many values to unpack (expected 2)}.
\end{enumerate}

\begin{lstlisting}[language=bash, caption={Mengukur akibat \textit{stage
  reordering}, skrip tersedia di
  \berkas{sources/chapter06/uji_susun_ulang.py}}, label=lst:jalankan-susun-ulang]
python uji_susun_ulang.py
# Keluarannya: Median urutan Resize dahulu    : 0.5835 detik
#              Median urutan Watermark dahulu : 0.6054 detik
#              Range urutan Resize            : 0.3660 detik
#              Range dibagi selisih           : 16.7 kali
#              Putaran Resize lebih cepat     : 88 dari 100
#              File sama persis               : False
\end{lstlisting}

Tabel~\ref{tab:lat2pipa-contoh} berisi hasil satu kali pengukuran nyata.

\begin{table}[htbp]
  \centering
  \caption{Contoh hasil Latihan 2 yang sudah terisi, seratus putaran}
  \label{tab:lat2pipa-contoh}
  \begin{tabularx}{\textwidth}{|X|l|}
    \hline
    \textbf{Yang dilaporkan} & \textbf{Nilai} \\
    \hline
    Median urutan \texttt{Resize} dahulu & 0,5835 detik \\
    \hline
    Median urutan \texttt{Watermark} dahulu & 0,6054 detik \\
    \hline
    Selisih median & 0,0219 detik \\
    \hline
    \textit{Range} dibagi selisih & 16,7 kali \\
    \hline
    Putaran \texttt{Resize} lebih cepat & 88 dari 100 \\
    \hline
    \textit{File} sama persis & tidak \\
    \hline
    \textit{Error} saat \texttt{SaveJpeg} dipindah & \texttt{ValueError} \\
    \hline
  \end{tabularx}
\end{table}

Isikan hasil pengukuran pada Tabel~\ref{tab:lat2pipa-hasil}.

\begin{table}[htbp]
  \centering
  \caption{Lembar isian Latihan 2}
  \label{tab:lat2pipa-hasil}
  \begin{tabularx}{\textwidth}{|X|l|}
    \hline
    \textbf{Yang dilaporkan} & \textbf{Nilai} \\
    \hline
    Median urutan \texttt{Resize} dahulu & \dots \\
    \hline
    Median urutan \texttt{Watermark} dahulu & \dots \\
    \hline
    \textit{Range} dibagi selisih & \dots \\
    \hline
    Putaran \texttt{Resize} lebih cepat & \dots \\
    \hline
    \textit{Error} saat \texttt{SaveJpeg} dipindah & \dots \\
    \hline
  \end{tabularx}
\end{table}

Jawab tiga pertanyaan berikut. Pertama, bandingkan \textit{range} dengan
selisih mediannya, lalu nyatakan apakah selisih waktu itu layak disebut ada.
Kedua, jelaskan mengapa \textit{file} keluarannya berbeda padahal
\textit{stage}-nya sama. Ketiga, jelaskan mengapa memindahkan \texttt{SaveJpeg}
justru gagal, lalu sebutkan apa yang sebenarnya mengikat kedua \textit{stage}
itu.

\subsection*{Latihan 3: \textit{First Output Latency}}

Latihan ini mengukur selisih \textit{first output latency} antara jalur
\textit{streaming} dan jalur \textit{batch}. Pengukurannya diulang seratus
putaran, sebab satu putaran tidak cukup untuk menyimpulkan apa pun. Langkah
kerjanya sebagai berikut.

\begin{enumerate}
  \item Jalankan perintah pertama pada Listing~\ref{lst:jalankan-ukur-aliran},
    lalu catat kesembilan angka pada laporan akhirnya.
  \item Buka \berkas{sources/chapter06/pipa.py}, lalu bandingkan isi
    \texttt{jalankan} dengan \texttt{jalankan\_sekaligus}. Tunjukkan satu kata
    yang membedakan keduanya.
  \item Jalankan kedua perintah \texttt{memori} pada listing yang sama, lalu
    catat puncak RSS masing-masing jalur.
  \item Turunkan \texttt{BATAS\_MEMORI} dari 60 menjadi 10, jalankan ulang kedua
    perintah \texttt{memori} itu, lalu catat perubahan angkanya. Satu kali
    percobaan nyata turun menjadi 40,0 MiB pada jalur \textit{streaming} dan
    118,1 MiB pada jalur \textit{batch}. Kembalikan nilainya sesudah dicatat.
\end{enumerate}

\begin{lstlisting}[language=bash, caption={Mengukur \textit{first output
  latency} beserta puncak RSS, skrip tersedia di
  \berkas{sources/chapter06/ukur_aliran.py}}, label=lst:jalankan-ukur-aliran]
python ukur_aliran.py
# Keluarannya: Median file pertama streaming : 83.1 ms
#              Median file pertama batch     : 792.8 ms
#              Putaran streaming lebih cepat : 100 dari 100
#              Puncak RSS streaming          : 118.4 MiB
#              Puncak RSS batch              : 576.8 MiB
python ukur_aliran.py memori streaming
# Keluarannya: 118.4
python ukur_aliran.py memori batch
# Keluarannya: 576.8
\end{lstlisting}

Tabel~\ref{tab:lat3pipa-contoh} berisi hasil satu kali pengukuran nyata.

\begin{table}[htbp]
  \centering
  \caption{Contoh hasil Latihan 3 yang sudah terisi, seratus putaran}
  \label{tab:lat3pipa-contoh}
  \begin{tabularx}{\textwidth}{|X|l|}
    \hline
    \textbf{Yang dilaporkan} & \textbf{Nilai} \\
    \hline
    Median \textit{file} pertama \textit{streaming} & 83,1 ms \\
    \hline
    Median \textit{file} pertama \textit{batch} & 792,8 ms \\
    \hline
    \textit{Batch} lebih lambat & 9,5 kali \\
    \hline
    \textit{Range} dibagi selisih & 0,18 kali \\
    \hline
    Putaran \textit{streaming} lebih cepat & 100 dari 100 \\
    \hline
    Puncak RSS \textit{streaming} & 118,4 MiB \\
    \hline
    Puncak RSS \textit{batch} & 576,8 MiB \\
    \hline
  \end{tabularx}
\end{table}

Isikan hasil pengukuran pada Tabel~\ref{tab:lat3pipa-hasil}.

\begin{table}[htbp]
  \centering
  \caption{Lembar isian Latihan 3}
  \label{tab:lat3pipa-hasil}
  \begin{tabularx}{\textwidth}{|X|l|l|}
    \hline
    \textbf{Yang dilaporkan} & \textbf{Enam puluh \textit{file}} & \textbf{Sepuluh \textit{file}} \\
    \hline
    Median \textit{file} pertama \textit{streaming} & \dots & \dots \\
    \hline
    Median \textit{file} pertama \textit{batch} & \dots & \dots \\
    \hline
    \textit{Range} dibagi selisih & \dots & \dots \\
    \hline
    Puncak RSS \textit{streaming} & \dots & \dots \\
    \hline
    Puncak RSS \textit{batch} & \dots & \dots \\
    \hline
  \end{tabularx}
\end{table}

Jawab tiga pertanyaan berikut dengan menunjuk angka pada tabel. Pertama,
jelaskan mengapa selisih pada latihan ini layak disebut ada, sedangkan selisih
pada Latihan~2 tidak. Kedua, sebutkan satu jenis \textit{stage} yang memaksa
\textit{pipeline} menunggu seluruh data, lalu jelaskan akibatnya bagi
\textit{file} pertama. Ketiga, dengan menggabungkan hasil ketiga latihan,
tentukan susunan mana yang akan dipakai bila jumlah \textit{file} naik menjadi
lima puluh ribu, lalu sebutkan angka yang menjadi dasar keputusannya.
